\documentclass[tikz,border=6pt]{standalone}
% Shared visual language for the thesis architecture figures.
% Compile each standalone source with XeLaTeX from this directory.
\usepackage[UTF8,scheme=plain,fontset=none]{ctex}
\usepackage{amsmath,amssymb}
\usepackage{fontspec}
\setmainfont{TeX Gyre Termes}
\setsansfont{TeX Gyre Heros}
\setmonofont{DejaVu Sans Mono}
\setCJKmainfont{Noto Serif SC}
\setCJKsansfont[BoldFont=SimHei]{Noto Sans SC}
\usetikzlibrary{arrows.meta,positioning,calc,fit,backgrounds,shapes.geometric,matrix,decorations.pathreplacing}

\definecolor{Ink}{HTML}{203040}
\definecolor{Muted}{HTML}{66788A}
\definecolor{LineSoft}{HTML}{AAB7C4}
\definecolor{DataBlue}{HTML}{2563A6}
\definecolor{DataFill}{HTML}{EAF3FB}
\definecolor{ClockPurple}{HTML}{7450A5}
\definecolor{ClockFill}{HTML}{F1ECF8}
\definecolor{ControlTeal}{HTML}{16827B}
\definecolor{ControlFill}{HTML}{E9F6F4}
\definecolor{DspPurple}{HTML}{6A4C93}
\definecolor{DspFill}{HTML}{F0EAF7}
\definecolor{RamGreen}{HTML}{367C59}
\definecolor{RamFill}{HTML}{ECF6F0}
\definecolor{FabricOrange}{HTML}{C56A1A}
\definecolor{FabricFill}{HTML}{FFF1E5}
\definecolor{EvidenceGreen}{HTML}{2F7D4B}
\definecolor{EvidenceFill}{HTML}{EAF6EE}
\definecolor{Amber}{HTML}{B07A16}
\definecolor{AmberFill}{HTML}{FFF6DD}
\definecolor{PendingGray}{HTML}{7C8792}
\definecolor{PendingFill}{HTML}{F1F3F5}
\definecolor{Danger}{HTML}{A04444}
\definecolor{DangerFill}{HTML}{FFF0F0}

\tikzset{
  >={Stealth[length=2.2mm,width=1.45mm]},
  every node/.style={font=\sffamily\footnotesize,text=Ink},
  data/.style={-{Stealth[length=2.2mm,width=1.45mm]},draw=DataBlue,line width=0.92pt},
  data bi/.style={{Stealth[length=1.9mm,width=1.3mm]}-{Stealth[length=1.9mm,width=1.3mm]},draw=DataBlue,line width=0.82pt},
  control/.style={-{Stealth[length=2.0mm,width=1.35mm]},draw=ControlTeal,dashed,line width=0.78pt},
  clock/.style={-{Stealth[length=2.0mm,width=1.35mm]},draw=ClockPurple,densely dotted,line width=0.9pt},
  address/.style={-{Stealth[length=2.0mm,width=1.35mm]},draw=FabricOrange,dash dot,line width=0.76pt},
  feedback/.style={-{Stealth[length=1.8mm,width=1.25mm]},draw=Muted,line width=0.66pt},
  block/.style={draw=Ink,fill=white,rounded corners=1.5pt,line width=0.68pt,minimum height=10mm,minimum width=23mm,align=center,inner sep=3pt},
  data block/.style={block,draw=DataBlue,fill=DataFill},
  control block/.style={block,draw=ControlTeal,fill=ControlFill},
  clock block/.style={block,draw=ClockPurple,fill=ClockFill},
  dsp/.style={block,draw=DspPurple,fill=DspFill,line width=0.9pt},
  ram/.style={block,draw=RamGreen,fill=RamFill,line width=0.9pt},
  fabric/.style={block,draw=FabricOrange,fill=FabricFill,line width=0.9pt},
  result/.style={block,draw=EvidenceGreen,fill=EvidenceFill,line width=0.9pt},
  pending/.style={block,draw=PendingGray,fill=PendingFill,line width=0.8pt},
  warning/.style={block,draw=Amber,fill=AmberFill,line width=0.9pt},
  boundary/.style={draw=Muted!75,rounded corners=2.2mm,dashed,line width=0.55pt,inner sep=3mm},
  title/.style={font=\sffamily\bfseries\large,text=Ink,align=center},
  subtitle/.style={font=\sffamily\scriptsize,text=Muted,align=center},
  group title/.style={font=\sffamily\bfseries\footnotesize,text=Ink,anchor=west},
  signal/.style={font=\sffamily\scriptsize,text=Ink,fill=white,inner sep=1.2pt,align=center},
  note/.style={font=\sffamily\scriptsize,text=Muted,align=center},
  tiny note/.style={font=\sffamily\tiny,text=Muted,align=center},
  badge/.style={draw=EvidenceGreen,fill=EvidenceFill,rounded corners=1.8mm,line width=0.65pt,font=\sffamily\bfseries\scriptsize,text=EvidenceGreen,inner xsep=4pt,inner ysep=2pt,align=center},
  pending badge/.style={draw=PendingGray,fill=PendingFill,rounded corners=1.8mm,line width=0.65pt,font=\sffamily\bfseries\scriptsize,text=PendingGray,inner xsep=4pt,inner ysep=2pt,align=center},
  innovation/.style={circle,draw=Amber,fill=AmberFill,line width=0.8pt,minimum size=6.5mm,font=\sffamily\bfseries\scriptsize,text=Amber,inner sep=0pt},
  sum/.style={circle,draw=Ink,fill=white,line width=0.75pt,minimum size=6mm,inner sep=0pt},
  mux/.style={trapezium,trapezium left angle=72,trapezium right angle=108,draw=Ink,fill=white,line width=0.68pt,minimum height=11mm,minimum width=12mm,align=center,inner sep=2pt}
}

\newcommand{\bitrate}[2]{{\scriptsize #1 bit，#2}}
\newcommand{\passmark}{\textcolor{EvidenceGreen}{\bfseries 满足判据}}
\newcommand{\pendingmark}{\textcolor{PendingGray}{\bfseries 待验证}}


\begin{document}
\begin{tikzpicture}[x=1cm,y=1cm]

% Title.
\node[title] at (13.0,11.55) {双采样率128倍插值系统的板级时钟、控制与数据协同架构};
\node[subtitle] at (13.0,11.10)
  {XC7A35T-FGG484-2；单音频主时钟域；$1\times/4\times/8\times/128\times$可观测输出};

% Clock plane.
\node[clock block,minimum width=20mm] (clk20) at (1.20,9.30) {20 MHz\\板载晶振};
\node[clock block,minimum width=19mm] (ibuf) at (3.45,9.30) {IBUF\\[-1pt]BUFG};
\node[clock block,minimum width=37mm,minimum height=17mm] (mmcm) at (6.65,9.30)
  {双MMCM时钟发生器\\[-1pt]
   {\scriptsize 44.1 kHz族：5.644796 MHz}\\[-1pt]
   {\scriptsize 48 kHz族：6.144068 MHz}};
\node[clock block,minimum width=25mm] (bufgmux) at (10.35,9.30) {BUFGMUX\_CTRL\\无毛刺选择};
\node[clock block,minimum width=31mm] (audioclk) at (13.55,9.30)
  {$clk_{audio,128\times}$\\[-1pt]{\scriptsize 全数据通路唯一主时钟}};
\node[control block,minimum width=30mm] (cegen) at (17.20,9.30)
  {7 bit CE/相位发生器\\[-1pt]{\scriptsize $F_s,2F_s,4F_s,8F_s,\ldots,128F_s$}};
\draw[clock] (clk20) -- (ibuf);
\draw[clock] (ibuf) -- (mmcm);
\draw[clock] (mmcm) -- (bufgmux);
\draw[clock] (bufgmux) -- (audioclk);
\draw[clock] (audioclk) -- (cegen);

% Control plane.
\node[control block,minimum width=27mm] (keypad) at (1.45,6.70) {4$\times$4矩阵键盘\\[-1pt]{\scriptsize 20 MHz控制域}};
\node[control block,minimum width=31mm,minimum height=17mm] (modectrl) at (5.05,6.70)
  {模式/采样率族控制\\[-1pt]
   {\scriptsize family：44.1/48 kHz}\\[-1pt]
   {\scriptsize mode：1/4/8/128$\times$}};
\node[control block,minimum width=32mm,minimum height=17mm] (cdc) at (9.10,6.70)
  {原子CDC握手\\[-1pt]
   {\scriptsize toggle/ack bundled-data}\\[-1pt]
   {\scriptsize 静音/稳定等待保护}};
\node[control block,minimum width=30mm] (commit) at (13.20,6.70)
  {音频域模式提交\\[-1pt]{\scriptsize 边界更新、无毛刺切换}};
\node[control block,minimum width=36mm,minimum height=17mm] (guard) at (17.35,6.70)
  {采样率族切换保护\\[-1pt]
   {\scriptsize family请求 + MMCM锁定}\\[-1pt]
   {\scriptsize 输出family\_active/audio reset}};
\draw[control] (keypad) -- (modectrl);
\draw[control] (modectrl) -- node[signal,above] {2 bit mode} (cdc);
\draw[control] (cdc) -- node[signal,above] {mode\_audio} (commit);
\draw[control] (modectrl.north) -- ++(0,0.42) -| node[signal,pos=0.78,above] {family\_sel} (guard.north west);
\draw[control] (mmcm.south) -- ++(0,-0.34) -| node[signal,pos=0.78,above] {locked} (guard.north);
\draw[control] (guard.north) -- ++(0,0.36) -| node[signal,pos=0.76,above] {family\_active} (bufgmux.south);

% Data plane.
\node[ram,minimum width=28mm,minimum height=16mm] (rom) at (1.50,3.65)
  {双采样率测试音ROM\\[-1pt]{\scriptsize signed 24 bit；$F_s$更新}\\[-1pt]{\scriptsize $x\_in\_valid=1$；板级RAMB18E1\#4}};
\node[data block,minimum width=35mm,minimum height=22mm] (core) at (5.65,3.65)
  {128倍插值核心\\[-1pt]
   {\scriptsize 三级$2\times$ FIR + CIC$16$}\\[-1pt]
   {\scriptsize 24/20/20 bit；2 DSP + 3 RAMB18}};
\node[mux,minimum height=19mm,minimum width=16mm] (tapmux) at (10.20,3.65) {4选1\\数据MUX};
\node[data block,minimum width=35mm,minimum height=20mm] (encode) at (14.05,3.65)
  {DAC编码与钳位\\[-1pt]
   {\scriptsize signed 24 bit $\rightarrow$}\\[-1pt]
   {\scriptsize 8 bit offset-binary}};
\node[data block,minimum width=29mm,minimum height=17mm] (iob) at (18.05,3.65)
  {下降沿IOB寄存\\[-1pt]{\scriptsize $dac\_data[7:0]$}\\[-1pt]{\scriptsize mute时强制中码128}};
\node[result,minimum width=26mm,minimum height=17mm] (dac) at (21.75,3.65)
  {AD9708\\[-1pt]{\scriptsize DAC数据/时钟接口}};
\draw[data] (rom) -- node[signal,above] {$x[n]$，24 bit} (core);
\draw[data] (rom.south) -- ++(0,-0.60) -| node[signal,pos=0.80,below] {raw旁路} (tapmux.south);
\draw[data] (core) -- node[signal,above] {$4/8/128\times$抽头} (tapmux);
\draw[data] (tapmux) -- (encode);
\draw[data] (encode) -- (iob);
\draw[data] (iob) -- (dac);

% Observable tap labels and DAC clock.
\node[note,anchor=north,align=left] at (8.50,1.50)
  {$1\times:x_{raw}$，$4\times:y_4$，$8\times:y_8$，$128\times:y_{128}$；$y_2$仅为内部节点};
\node[clock block,minimum width=26mm] (oddr) at (18.05,1.60) {ODDR/分频控制\\[-1pt]{\scriptsize $dac\_clk$生成}};
\draw[clock] (audioclk.south) |- (oddr.west);
\draw[clock] (oddr.east) -| (dac.south);

% Cross-plane connections.
\draw[clock] (audioclk.south) |- (core.north);
\draw[control] (cegen.south) |- node[signal,pos=0.70,right] {CE/phase} (core.north east);
\draw[control] (commit.south) |- (tapmux.north);
\draw[control] (commit.east) -| node[signal,pos=0.36,right] {系数模式} (core.north east);
\draw[control] (cdc.south) -- ++(0,-0.35) -| node[signal,pos=0.84,below] {audio\_mute} (iob.north);
\draw[control] (guard.south) -- ++(0,-0.35) -| node[signal,pos=0.82,below] {audio reset} (core.north west);

% Innovation callouts.
\node[innovation] (i1) at (19.50,9.30) {I1};
\node[note,anchor=west,align=left] at (19.90,9.30) {单时钟+事件使能\\实现多速率};
\node[innovation] (i2) at (20.20,7.05) {I2};
\node[note,anchor=west,align=left] at (20.60,7.05) {原子CDC与静音保护\\约束异步模式切换};
\node[innovation] (i3) at (20.20,6.25) {I3};
\node[note,anchor=west,align=left] at (20.60,6.25) {MMCM锁定参与\\无毛刺采样率族切换};
\node[innovation] (i4) at (22.00,2.15) {I4};
\node[note,anchor=west,align=left] at (22.40,2.15) {单核心多倍率抽头\\加$1\times$ raw旁路};

% Background boundaries and legend.
\begin{scope}[on background layer]
  \node[boundary,fill=ClockFill!40,fit=(clk20)(ibuf)(mmcm)(bufgmux)(audioclk)(cegen)(i1)] (clockPlane) {};
  \node[boundary,fill=ControlFill!42,fit=(keypad)(modectrl)(cdc)(commit)(guard)(i2)(i3)] (ctrlPlane) {};
  \node[boundary,fill=DataFill!30,fit=(rom)(core)(tapmux)(encode)(iob)(dac)(oddr)(i4)] (dataPlane) {};
\end{scope}
\node[group title] at ($(clockPlane.north west)+(0.2,-0.25)$) {时钟平面};
\node[group title] at ($(ctrlPlane.north west)+(0.2,-0.25)$) {控制与跨域平面};
\node[group title] at ($(dataPlane.north west)+(0.2,-0.25)$) {音频数据平面};

\draw[data] (0.80,0.55) -- (2.05,0.55) node[signal,right] {样值数据};
\draw[control] (4.20,0.55) -- (5.45,0.55) node[signal,right] {CE/valid/模式控制};
\draw[clock] (8.65,0.55) -- (9.90,0.55) node[signal,right] {物理时钟};
\node[note,anchor=west,align=left] at (13.05,0.55)
  {注：$F_s$至$128F_s$表示数据有效事件率；三级FIR与CIC之间不存在数据时钟域跨越。};

\end{tikzpicture}
\end{document}


